\section{Problem Formulation}
\label{sec:formulation}

\subsection{Agentic Evolution}
\label{sec:agentic_evolution}

We formulate agentic evolution as a process for optimizing an artifact through repeated improvement rounds. 
Let $x \in \mathcal{X}$ denote the object being optimized, such as a program, prompt, workflow, skill, tool, or agent component. 
We use $r$ to index the evolution round. 
Each round produces a round context $c_r$, which contains the candidates generated in that round, their evaluation results, execution traces, failures, costs, and any intermediate information produced during optimization. 
The accumulated evolution context after $r$ rounds is denoted as
\[
\mathcal{C}_r = (c_1, c_2, \ldots, c_r).
\]
Finally, let $\Pi$ denote the optimization mechanism that advances evolution:
\[
c_r = \Pi(\mathcal{C}_{r-1}), 
\qquad 
\mathcal{C}_r = \mathcal{C}_{r-1} \oplus c_r,
\]
where $\oplus$ appends the newly produced round context to the accumulated evolution context. 
$\Pi$ does not have to be a fixed algorithm; it can also be an agentic process that reads the history, reasons over feedback, and decides how to generate the next candidate. 
Thus, $\Pi$ represents the mechanism by which search is continued from the current evolution context.

Under this formulation, existing agentic evolution methods mainly differ in how $\Pi$ is instantiated. 
In \textbf{procedure-based evolution }, $\Pi$ is a predefined outer loop whose behavior is mainly determined by selection and optimization: the selection rule chooses previous candidates or contexts from $\mathcal{C}_{r-1}$, while the optimization operator generates new candidates from the selected information. 
Evaluation assigns scores, traces, and feedback to the generated candidates, providing signals for future selection and update. 
In \textbf{agent-based evolution}, $\Pi$ is instead implemented by a general-purpose agent. 
Rather than following fixed selection-and-optimization rules, the agent reads the accumulated context $\mathcal{C}_{r-1}$ and decides what to do next, such as inspecting feedback, comparing candidates, modifying artifacts, writing tools, or generating new attempts. 
Thus, procedure-based evolution  specifies search control explicitly but rigidly, while agent-based evolution leaves search control implicit in the agent's context-conditioned behavior.

In both cases, evolution proceeds by repeatedly applying $\Pi$ while accumulating context $\mathcal{C}_r$. 
This context records not only the candidates produced by evolution, but also how search has unfolded through evaluation results, traces, failures, costs, and intermediate artifacts. 
The next subsection uses this accumulated context to define an environment view of evolution.

\subsection{Evolution as an Interactive Environment}
\label{sec:evolution_as_environment}

We treat the evolution process itself as an interactive environment for a meta-agent. 
At round $r$, the state of this environment is defined by the round index and the accumulated evolution context:
\[
s_r = (r, \mathcal{C}_r).
\]
When the optimization mechanism may change across rounds, we write the current mechanism as $\Pi_r$. 
This mechanism specifies the transition of the environment. 
Without intervention, the next round is produced by applying the current mechanism to the current context:
\[
c_{r+1} = \Pi_r(\mathcal{C}_r),
\qquad
s_{r+1} = (r+1, \mathcal{C}_r \oplus c_{r+1}).
\]
Thus, $\Pi_r$ is the transition rule that determines how the evolution process continues.

To interact with this environment, we introduce a meta-agent $M$. 
The role of $\Pi_r$ is to continue the candidate search, while the role of $M$ is to act on the evolution process that governs this search. 
Since the full state $s_r$ can be large and noisy, the meta-agent receives an observation extracted from the state:
\[
o_r = \Phi(s_r) = \Phi(r, \mathcal{C}_r),
\]
where $\Phi$ summarizes relevant information from the accumulated context, such as progress, repeated failures, invalid attempts, cost patterns, or redundant search directions.

Given this observation, the meta-agent produces an edit action:
\[
a_r = M(o_r).
\]
The action does not directly become the next candidate. 
Instead, it modifies the transition rule of the evolution environment:
\[
\Pi_{r+1} = \mathrm{Edit}(\Pi_r, a_r).
\]
The edited mechanism is then used to continue evolution:
\[
c_{r+1} = \Pi_{r+1}(\mathcal{C}_r),
\qquad
\mathcal{C}_{r+1} = \mathcal{C}_r \oplus c_{r+1}.
\]
In this sense, we formulate agentic evolution as an environment in which the state is the accumulated evolution context, the observation is a summary of this context, and the action edits the mechanism that controls future search.

This formulation applies to both forms of agentic evolution. 
For procedure-based evolution , editing $\Pi_r$ changes explicit components such as selection, optimization, feedback use, budget allocation, or update rules. 
For agent-based evolution, editing $\Pi_r$ changes the agentic context that shapes future decisions, such as skills, goals, tools, feedback format, or execution context. 
In both cases, the meta-agent steers evolution not by proposing one more candidate, but by modifying how subsequent search is carried out. 
Section~\ref{sec:methodology} describes the system design used to instantiate this formulation.
